\documentclass[11pt]{article}
\usepackage[margin=25mm]{geometry}
\usepackage{times}
\usepackage{hyperref}
\pagestyle{empty}

\begin{document}

\begin{center}
{\large\bfseries Cover Letter}
\end{center}

\vspace{0.5cm}

\noindent\textbf{To:} The Editor-in-Chief\\
\noindent\textbf{Journal:} \emph{Remote Sensing} (MDPI)\\
\noindent\textbf{Date:} October 16, 2026

\vspace{0.5cm}

\noindent Dear Editor,

We are pleased to submit our manuscript entitled ``The Split Is Part of the
Measurement: Spatial-Interleaving Evaluation Bias in Change Detection,
Quantified in Bi-Temporal Mangrove Mapping from Sentinel-2'' for
consideration for publication in \emph{Remote Sensing}.

This paper documents a spatial-evaluation failure mode in bi-temporal change
detection: a spatially interleaved random split produces a 29.6-percentage-point
higher change $F1$ than a buffer-separated test assignment, positive in all ten
partitions, while the corresponding ROC-AUC difference is only $+1.9$\,pp. The
audit quantifies this in a Sentinel-2 mangrove change-detection setting
(0.8--2.6\% change prevalence), decomposes the gap into a small
prevalence-arithmetic component ($+2.7$\,pp) and a dominant recall shift
associated with the joint spatial configuration of training and test samples,
and states the identification boundary of what the design can and cannot
attribute.

We believe this work is well suited to \emph{Remote Sensing} because it
addresses a methodological concern that affects a large share of change-detection
studies: the evaluation protocol itself moves the reported score. The paper
provides a quantified, reproducible demonstration of the effect size, a
decomposition that separates prevalence arithmetic from the spatial-configuration
component, and a ten-item reporting checklist for practitioners.

The manuscript is original, has not been published previously, and is not under
consideration elsewhere. All authors have read and approved the submitted
version. The full code base and results are deposited at
\url{https://github.com/hao-feier/split-audit}.

Thank you for your consideration.

\vspace{1cm}

\noindent Sincerely,\\
Junzhe Yuan\\
School of Environment and Ecology, Hunan Agricultural University, Changsha, China\\
\texttt{haofeier0613@gmail.com}\\[4pt]
Feier Hao\\
College of Information and Intelligence, Hunan Agricultural University, Changsha, China

\end{document}
